在 GLWE 密码体制中，模数切换是将密文的模数域改为更小（或更大）的模数域，同时保证密文仍能解密得到相同明文的过程。例如，设有密文 $\textsf{LWE}_{S, \sigma}(\Delta m) \in \mathbb{Z}_q^{k+1}$。如果将密文模数按 $q \rightarrow q'$ 切换，则密文变为 $\textsf{LWE}_{S, \sigma}\left(\Delta\dfrac{q'}{q} m + e'\right) \in \mathbb{Z}_{q'}^{k+1}$。密文的其他分量，例如噪声（$e$）和公钥分量（$a_0, a_1, \cdots, a_{k-1}$、$b$），均按 $\dfrac{q'}{q}$ 缩放，变为 $\left\lceil e\dfrac{q'}{q} \right\rfloor$ 和 $\left(\left\lceil a_0\dfrac{q'}{q}\right\rfloor, \left\lceil a_1\dfrac{q'}{q}\right\rfloor, \cdots, \left\lceil a_{k-1}\dfrac{q'}{q}\right\rfloor, \left\lceil b\dfrac{q'}{q}\right\rfloor\right)$。为切换 LWE 密文的模数，我们使用 \autoref{sec:modulus-rescaling} 介绍的模数重缩放技术。
相同的模数切换技术也适用于 RLWE 密文。本节将介绍如何切换（即重缩放）LWE 和 RLWE 密文的模数，并证明其正确性。
%That is, after  LWE/RLWE ciphertexts, we will prove why they are still valid LWE/RLWE ciphertexts.


\subsection{LWE 模数切换}
\label{subsec:modulus-switch-lwe}

\textbf{- 参考资料：}
\href{https://www.jeremykun.com/2022/07/16/modulus-switching-in-lwe/}{LWE 中的模数切换}

$ $

回顾 LWE 密码体制（\autoref{subsec:lwe-enc}），其组成如下：

\begin{itemize}
\item \textbf{\underline{参数设置}：} $\Delta = \dfrac{q}{t}$，\text{ } $S = (s_0, s_1, \ldots, s_{k-1}) \xleftarrow{\$} \mathbb{Z}_2^{k}$

\item \textbf{\underline{加密输入}：} $m \in \mathbb{Z}_t$，\text{ } $\vv{a} = (a_0, a_1, \ldots, a_{k-1}) \xleftarrow{\$} \mathbb{Z}_q^{k}$，\text{ } $e \xleftarrow{\chi_\sigma} \mathbb{Z}_q$

$ $

\item \textbf{\underline{加密}：} $\textsf{LWE}_{\vv{s},\sigma}(\Delta  m + e) = (\vv{a}, b) \in \mathbb{Z}_q^{k + 1}$ \text{ }（其中 $b = \vv{a} \cdot \vv{s} + \Delta  m + e \in \mathbb{Z}_q$）

$ $

\item \textbf{\underline{解密}：} $\textsf{LWE}^{-1}_{S,\sigma}(\textsf{ct}) = b - \vv{a}\cdot \vv{s} = \Delta  m + e \text{ } \text{ } \in \mathbb{Z}_q$

\end{itemize}

$ $


在 LWE 密码体制中，模数切换将原 LWE 密文的模数域转换成较小的模数域。可以将这一过程看作：除明文 $m$ 和私钥 $S$ 外，将原密文的所有分量缩小到更小的域中。该操作保持原明文 $m$ 的数值和完整性，同时将尺度因子 $\Delta$ 缩小为 $\hat{\Delta}$，将噪声 $e$ 缩小为 $\hat{e}$（注意，噪声的变化不影响原明文 $m$，因为噪声最终会在解密后的舍入过程中被消去），并将 $\vv{a}$ 缩小为 $\vv{\hat{a}}$。TFHE 在自举过程中使用模数切换来提高计算效率（详见 \autoref{subsec:tfhe-noise-bootstrapping}）。此外，BGV（\autoref{subsec:bfv-mult-cipher}）和 CKKS（\autoref{subsec:ckks-mult-cipher}）的密文与密文乘法算法也会使用模数切换。

LWE 模数切换的核心思想是对 LWE 方案中的同余关系进行重缩放。LWE 的同态运算包括密文与密文加法、密文与明文加法、密文与明文乘法以及密文与密文乘法。这些算法中使用的同余关系本质上都可由同一个基本关系改写得到：$b = \vv{a}\cdot \vv{s} + \Delta m + e \bmod q$。因此，将 LWE 密文的模数按 $q \rightarrow q'$ 切换，等价于将上述同余关系的模数按 $q \rightarrow q'$ 重缩放。

由此可知，在 LWE 密码体制中按 $q \rightarrow \hat{q}$（其中 $q > \hat{q}$）进行模数切换，就是按如下方式转换原 LWE 密文 $\textsf{LWE}_{S, \sigma}(\Delta m + e)$：

\begin{tcolorbox}[title={\textbf{\tboxlabel{\ref*{subsec:modulus-switch-lwe}} LWE 模数切换}}]

给定满足 $b = \vv{a}\cdot \vv{s} + \Delta m + e \bmod q$ 且 $m \in \mathbb{Z}_t$ 的 LWE 密文 $(\vv{a}, b)$，将其模数从 $q$ 切换到 $\hat q$，等价于将 $(A, b)$ 更新为 $(\vv{\hat{a}}, \hat b)$，具体如下：


$\vv{\hat{a}} = (\hat{a}_0, \hat{a}_1, \ldots, \hat{a}_{k-1})$，其中各分量满足 $\hat{a}_i = \Big\lceil a\dfrac{\hat{q}}{q} \Big\rfloor \in \mathbb{Z}_{\hat{q}}$ \textcolor{red}{ $\rhd$ $\lceil \rfloor$ 表示舍入到最近的整数}

 $\hat{b} = \Big\lceil b \dfrac{\hat{q}}{q} \Big\rfloor \in \mathbb{Z}_{\hat{q}}$

$\textsf{LWE}_{{S},\sigma}(\hat{\Delta}  m + \hat{e} + \epsilon_{\textit{all}}) = (\hat{a}_0, \hat{a}_1, \ldots, \hat{b}) \in \mathbb{Z}_{\hat{q}}^{k + 1}$ \textcolor{red}{ $\rhd$ 其中 $\epsilon_{\textit{all}}$ 是一个较小的舍入误差}

$ $


上述更新使 $\hat e$ 和 $\hat \Delta$ 发生如下变化：

$\hat{e} = \Big\lceil e\dfrac{\hat{q}}{q} \Big\rfloor \in \mathbb{Z}_{\hat{q}}$, \text{ }

$\hat{\Delta} = \Delta\dfrac{\hat{q}}{q}$ \textcolor{red}{ $\rhd$ 该值应为整数}

$ $


与此同时，$S$ 和 $m$ 保持不变。

\end{tcolorbox}

$ $

注意，要说明 $(\hat{a}_0, \hat{a}_1, \ldots, \hat{b}) \in \mathbb{Z}_{\hat{q}}^{k + 1}$ 是 $\hat{\Delta}m$ 的有效 LWE 密文，需要证明以下关系成立：

$\hat{b} = \sum\limits_{i=0}^{k-1}\hat{a}_i \cdot s_i + \hat{\Delta} m + \hat{e} \in \mathbb{Z}_{\hat{q}}$

$ $

%\textbf{\underline{Proof}}
\begin{myproof}
\begin{enumerate}
\item 注意以下关系：

$\hat{b} = \Big\lceil b \dfrac{\hat{q}}{q} \Big\rfloor = b\dfrac{\hat{q}}{q} + \epsilon_b$（其中 $-0.5 < \epsilon_b < 0.5$，表示舍入产生的偏差）

$\hat{a_i} = \Big\lceil a_i \dfrac{\hat{q}}{q} \Big\rfloor = a_i\dfrac{\hat{q}}{q} + \epsilon_{a_i}$（其中 $-0.5 < \epsilon_{a_i} < 0.5$）

$\hat{e} = \Big\lceil e \dfrac{\hat{q}}{q} \Big\rfloor = e\dfrac{\hat{q}}{q} + \epsilon_e$（其中 $-0.5 < \epsilon_e < 0.5$）
\item 注意以下关系：

$b = \vv{a} \cdot \vv{s} + \Delta  m + e = \sum\limits_{i=0}^{k-1}(a_is_i) + \Delta m + e  \in \mathbb{Z}_q$

$b = \sum\limits_{i=0}^{k-1}(a_is_i) + \Delta m + e + H \cdot q$（这里将模 $q$ 的关系改写为加上 $q$ 的某个倍数 $H \cdot q$）

\item 根据步骤 1 和步骤 2：

$\hat{b} = b\dfrac{\hat{q}}{q} + \epsilon_b \in \mathbb{Z}_{\hat{q}}$

$= \left(\sum\limits_{i=0}^{k-1}(a_is_i) + \Delta m + e + H \cdot q\right) \cdot \dfrac{\hat{q}}{q}  + \epsilon_b$

$= \dfrac{\hat{q}}{q} \cdot \sum\limits_{i=0}^{k-1}(a_is_i) + \dfrac{\hat{q}}{q} \cdot \Delta m + \dfrac{\hat{q}}{q} \cdot e + \dfrac{\hat{q}}{q} \cdot H \cdot q + \epsilon_b$

$= \sum\limits_{i=0}^{k-1}\left(\dfrac{\hat{q}}{q} \cdot a_is_i\right) + \hat{\Delta} m + (\hat{e} - \epsilon_e) + \hat{q}\cdot H + \epsilon_b$

$= \sum\limits_{i=0}^{k-1}\left((\hat{a}_i - \epsilon_{a_i}) \cdot s_i\right) + \hat{\Delta} m + \hat{e} - \epsilon_e + \hat{q}\cdot H + \epsilon_b$

$= \sum\limits_{i=0}^{k-1}(\hat{a}_is_i - \epsilon_{a_i}s_i) + \hat{\Delta} m + \hat{e} - \epsilon_e + \epsilon_b \in \mathbb{Z}_{\hat{q}}$

$= \sum\limits_{i=0}^{k-1}\hat{a}_is_i + \hat{\Delta} m + \left( \hat{e} - \epsilon_e + \epsilon_b - \sum\limits_{i=0}^{k-1}\epsilon_{a_i}s_i \right) \in \mathbb{Z}_{\hat{q}}$

$= \sum\limits_{i=0}^{k-1}\hat{a}_is_i + \hat{\Delta} m + \hat{e} + \epsilon_{\textit{all}} \in \mathbb{Z}_{\hat{q}}$ \textcolor{red}{ $\rhd$ 其中 $\epsilon_{\textit{all}} = \left(- \epsilon_e + \epsilon_b - \sum\limits_{i=0}^{k-1}\epsilon_{a_i}s_i \right)$}

$ $

$\epsilon_{\textit{all}}$ 的最大可能值为：

$\epsilon_{\textit{all}} = |-0.5| + |0.5| + |-0.5 \cdot k| = 1 + 0.5k$ 

因此，LWE 模数切换得到的是一个近似同余关系（\autoref{sec:modulus-rescaling}）。不过，只要 $\hat \Delta$ 足够大，$\epsilon_{\textit{all}} = 1 + 0.5k$ 就会在 LWE 解密的右移过程中被消去，最终恢复原明文 $m$。此外，在实际应用中，当 $k$ 足够大时，$\sum\limits_{i=0}^{k-1}\epsilon_{a_i}s_i$
相对于密文模数仍较小，因为每个 $a_i$ 和 $s_i$ 都由均匀采样产生。

\para{注意：} 如果 $\hat \Delta$ 不够大，$\epsilon_{\textit{all}}$ 就可能无法在解密时被消去，从而破坏明文 $m$。此外，如果 $\Delta \rightarrow \hat\Delta$ 缩小得过多，$\hat\Delta m$ 与 $\hat e$ 之间的距离会变得过小；此时 $\hat e = \Big\lceil e \dfrac{\hat{q}}{q} \Big\rfloor$ 的舍入过程可能影响 $\hat \Delta m$ 的最低有效位，从而破坏明文。

$ $

\item 综上，$\hat{b}$ 近似满足：

$\hat{b} = \sum\limits_{i=0}^{k-1}\hat{a}_is_i + \hat{\Delta} m + \hat{e} + \epsilon_{\textit{all}}  \approx \sum\limits_{i=0}^{k-1}\hat{a}_i  s_i + \hat{\Delta} m + \hat{e} \in \mathbb{Z}_{\hat{q}}$

因此，$(\hat{a}_0, \hat{a}_1, \ldots, \hat{b}) = \textsf{LWE}_{{S},\sigma}(\hat{\Delta} m + \hat{e} + \epsilon_{\textit{all}})$，对其解密即可得到 $m$。

%\begin{flushright}
%\qedsymbol{} 
%\end{flushright}

\end{enumerate}
\end{myproof}

\subsection{示例}

假设 LWE 的参数设置如下：

$ $

$t = 4$

$q = 64$

$n = 4$

$\Delta = \dfrac{q}{t} = 16$

$m = 1 \in \mathbb{Z}_t$

$S = (s_0, s_1, s_2, s_3) = (0, 1, 1, 0) \in \{-1, 0, 1\}^4$

$A = (a_0, a_1, a_2, a_3) = (-25, 12, -3, 7) \in \mathbb{Z}_q^4$

$e = 1 \in \mathbb{Z}_q$

$b = a_0s_0 + a_1s_1 + a_2s_2 + a_3s_3 + \Delta m + e = 26 \in \mathbb{Z}_q$

$\textsf{LWE}_{S, \sigma}(\Delta m + e) = \textsf{ct} = (a_0, a_1, a_2, a_3, b) = (-25, 12, -3, 7, 26) \in \mathbb{Z}_q^{n+1}$

$ $

现在希望将模数从 $q=64$ 切换到 $\hat{q} = 32$，可得：

$\hat{\Delta} = \Delta\cdot\dfrac{32}{64} = 8$

$\hat{e} = \Big\lceil 1 \cdot \dfrac{32}{64} \Big\rfloor = 1$

$\textsf{LWE}_{S, \sigma}(\hat{\Delta} m + \hat{e} + \epsilon_{\textit{all}}) = \hat{\textsf{ct}} = (\hat{a_0}, \hat{a_1}, \hat{a_2}, \hat{a_3}, \hat{b})$

$ = \left(\Big\lceil -25\cdot\dfrac{32}{64}\Big\rfloor, \Big\lceil12\cdot\dfrac{32}{64}\Big\rfloor, \Big\lceil-3\cdot\dfrac{32}{64}\Big\rfloor, \Big\lceil7\cdot\dfrac{32}{64}\Big\rfloor, \Big\lceil26\cdot\dfrac{32}{64}\Big\rfloor\right)$

$ = (-12, 6, -1, 4, 13) \in \mathbb{Z}_{\hat{q}}^{n+1}$

$ $

接下来验证以下 LWE 约束是否成立：

$\hat{b} = \hat{a}_0s_0 + \hat{a}_1s_1 + \hat{a}_2s_2 + \hat{a}_3s_3 + \hat{\Delta}m + \hat{e} \in \mathbb{Z}_{32}$

$13 = 0 + 6 - 1 + 0 + 8 \cdot 1 + 1 \in \mathbb{Z}_{32}$

$13 \approx 14 \in \mathbb{Z}_{32}$

此处出现了大小为 1 的微小差异，来源于以下舍入偏差：

$\hat {a_0} = \lceil -12.5 \rfloor = -12$、$\hat {a_2} = \lceil -1.5 \rfloor = -1$、$\hat{a_3} =  \lceil 3.5 \rfloor = 4$ 和 $\hat{e} = \lceil 0.5 \rfloor = 1 $

$ $

代入 LWE 解密公式计算：

$\hat b - (\hat{a}_0s_0 + \hat{a}_1s_1 + \hat{a}_2s_2 + \hat{a}_3s_3) = 13 - (0(-12) + 1(6) + 1(-1) + 0(4)) = 13 - (6 - 1) = 8 = \hat m + \hat e \in \mathbb{Z}_{32}$

$ $

$m = \left \lceil \dfrac{8}{\hat \Delta} \right \rfloor = \left \lceil \dfrac{8}{8} \right \rfloor = 1$，结果正确。


\subsection{讨论}
\label{subsubsec:modulus-switch-lwe-discuss}

\begin{figure}[h!]
    \centering
  \includegraphics[width=0.7\linewidth]{figures/modulus-switching.pdf}
  \caption{含噪声的缩放明文示意图：$\Delta \cdot m + e \in \mathbb{Z}_q$}
  \label{fig:modulus-switch}
\end{figure}

\para{$m$ 与 $e$ 之间的距离缩小：} 将 LWE 密文的模数按 $q \rightarrow \hat{q}$ 切换后，底层含噪明文 $\Delta m + e$ 缩小为 $\hat{\Delta}m + \hat{e}$，如 \autoref{fig:modulus-switch} 所示。注意，在按 $q \rightarrow \hat{q}$ 切换模数后，$\Delta m$ 缩小为 $\hat{\Delta} m$，但不会丢失明文比特信息。具体而言，明文值 $m$ 保持不变，而尺度因子 $\Delta$ 缩小为 $\hat{\Delta}$，噪声 $e$ 缩小为 $\hat e$。不过，与 $e$ 的最高有效位和 $\Delta m$ 的最低有效位之间的距离相比，模数切换后 $\hat e$ 的最高有效位与 $\hat \Delta m$ 的最低有效位之间的距离更小。




%\para{Hard Threshold on the $\Delta m$ Term's Value:} In \autoref{subsubsec:glwe-add-cipher-discuss}, we learned that the value of $\Delta m$ should not overflow or underflow its MSB $q$ during homomorphic computation (i.e., $\Delta m$ should be within the range of modulo $q$ without wrapping around), because if it wraps around the MSB, then our derived congruence relation of modulus switch loses its correctness. This is necessary in order to get the correct functionality of modulus switch on LWE ciphertexts. In particular, if we need to handle the case of $\Delta m$'s overflow, then our congruence relation of modulus switch should include the additional term $H_2 \cdot \Delta$ to take into account the added multiples of $\Delta$ caused by the overflows of $\Delta m$. But problematically, this term does not get eliminated during the modulus switch as $H_2 \cdot \Delta$ is not divisible by $\hat{q}$, which breaks our derived congruence relation above. Therefore, the applications of GLWE homomorphic encryption should ensure to manage their scaled plaintext $\Delta m$ to be always within than its maximum upper and minimum lower boundaries of the ciphertext modulus domain (i.e., $[0, q - 1]$ or $[-\dfrac{q}{2}, \dfrac{q}{2} - 1]$) throughout all computations. Meanwhile, the masking terms $A_i\cdot S_i$ are allowed to overflow or underflow ciphertext indefinitely during homomorphic computations, because during the decryption procedure they will be exactly subtracted and become zeroed out modulo $q$.  


\subsection{RLWE 模数切换}
\label{subsec:modulus-switch-rlwe}




RLWE 模数切换与 LWE 模数切换类似。回顾 RLWE 密码体制（\autoref{subsec:rlwe-enc}），其组成如下：

\begin{itemize}
\item \textbf{\underline{参数设置}：}
$\Delta = \dfrac{q}{t}$, \text{ } $S = s_0 + s_1X + s_2X^2 + \cdots + s_{n-1}X^{n-1} \xleftarrow{\$} \mathcal{R}_{\langle n, \textit{tern} \rangle}$

$ $

\item \textbf{\underline{加密输入}：}

$M \in \mathcal{R}_{\langle n, t \rangle}$

$A = a_0 + a_1X + a_2X^2 + \cdots + a_{n-1}X^{n-1} \xleftarrow{\$} \mathcal{R}_{\langle n, q \rangle}$

$E = e_0 + e_1X + e_2X^2 + \cdots + e_{n-1}X^{n-1} \xleftarrow{\chi_\sigma} \mathcal{R}_{\langle n, q \rangle}$ 

$ $

\item \textbf{\underline{加密}：}
$\textsf{RLWE}_{S,\sigma}(\Delta \cdot M + E) = (A, B) \in \mathcal{R}_{\langle n,q \rangle}^2$ 

其中 $B = A \cdot S + \Delta \cdot M + E = b_0 + b_1X + b_2X^2 + \cdots + b_{n-1}X^{n-1}$

$ $

\item \textbf{\underline{解密}：} $\textsf{RLWE}^{-1}_{S,\sigma}(\textsf{ct}) = B - A \cdot S = \Delta  M + E \text{ } \in \mathcal{R}_{\langle n,q \rangle}$
\end{itemize}

$ $

RLWE 模数切换过程如下：

\begin{tcolorbox}[title={\textbf{\tboxlabel{\ref*{subsec:modulus-switch-rlwe}} RLWE 模数切换}}]

对于满足 $B = AS + \Delta M + E$ 且 $M \in \mathcal{R}_{\langle n, q \rangle}$ 的 RLWE 密文 $(A, B)$，将模数从 $q$ 切换到 $\hat q$，等价于将 $(A, B)$ 更新为 $(\hat A, \hat B)$，具体如下：


$\hat{A} = \hat{a}_0 + \hat{a}_1X + \hat{a}_2X^2 + \cdots + \hat{a}_{n-1}X^{n-1}$，其中每个系数满足 $\hat{a}_i = \Big\lceil a_i\dfrac{\hat{q}}{q} \Big\rfloor \in \mathbb{Z}_{\hat{q}}$

$\hat{B} = \hat{b}_0 + \hat{b}_1X + \hat{b}_2X^2 + \cdots + \hat{b}_{n-1}X^{n-1}$，其中每个系数满足 $\hat{b}_i = \Big\lceil b_i\dfrac{\hat{q}}{q} \Big\rfloor \in \mathbb{Z}_{\hat{q}}$

$\textsf{RLWE}_{{S},\sigma}(\hat{\Delta}  M + \hat{E} + E^{\langle \epsilon_{\textit{all}} \rangle}) = (\hat{A}, \hat{B}) \in \mathcal{R}_{\langle n, \hat{q} \rangle}^{2}$ 

$ $

上述更新使 $\Delta$ 和 $E$ 发生如下变化：

$\hat{\Delta} = \Delta\dfrac{\hat{q}}{q}$ \textcolor{red}{ $\rhd$ 该值应为整数}

$\hat{E} = \hat{e}_0 + \hat{e}_1X + \hat{e}_2X^2 + \cdots + \hat{e}_{n-1}X^{n-1}$，其中每个系数满足 $\hat{e}_i = \Big\lceil e_i\dfrac{\hat{q}}{q} \Big\rfloor \in \mathbb{Z}_{\hat{q}}$

$ $

与此同时，$S$ 和 $M$ 保持不变。

\end{tcolorbox}

证明与 LWE 模数切换类似。

$ $

\textbf{\underline{证明}}

\begin{enumerate}
\item 注意以下关系：

$\hat{b}_i = \Big\lceil b_i \dfrac{\hat{q}}{q} \Big\rfloor = b_i\dfrac{\hat{q}}{q} + \epsilon_{b_i}$（其中 $-0.5 < \epsilon_{b_i} < 0.5$，表示舍入产生的偏差）

$\hat{a}_i = \Big\lceil a_i \dfrac{\hat{q}}{q} \Big\rfloor = a_i\dfrac{\hat{q}}{q} + \epsilon_{a_i}$（其中 $-0.5 < \epsilon_{a_i} < 0.5$）

$\hat{e}_i = \Big\lceil e_i \dfrac{\hat{q}}{q} \Big\rfloor = e_i\dfrac{\hat{q}}{q} + \epsilon_{e_i}$（其中 $-0.5 < \epsilon_{e_i} < 0.5$）

\item 注意以下关系：


$B - A\cdot S$ 

$ = (b_0 + b_1X + \cdots + b_{n-1}X^{n-1} ) - (a_{0} + a_{1}X + \cdots + a_{n-1}X^{n-1})(s_{0} + s_{1}X + \cdots + s_{n-1}X^{n-1})$ 


$ = \left(b_0 - \left( \sum\limits_{i=0}^{0}(a_{0-i}s_{i}) - \sum\limits_{i=1}^{n-1}(a_{n-i}s_{i}) \right)\right)$

$ + \left(b_1 - \left( \sum\limits_{i=0}^{1}(a_{1-i}s_{i}) - \sum\limits_{i=2}^{n-1}(a_{n+1-i}s_{i})   \right) \right)\cdot X$ 



$ + \left(b_2 - \left( \sum\limits_{i=0}^{2}(a_{2-i}s_{i}) - \sum\limits_{i=3}^{n-1}(a_{n+2-i}s_{i})   \right) \right)\cdot X^2$ 


$\vdots$ 


$ + \left(b_{n-1} - \left(  \sum\limits_{i=0}^{n-1}(a_{n-1-i}s_{i}) -  \sum\limits_{i=n}^{n-1}(a_{n+n-1-i}s_{i})  \right) \right)\cdot X^{n-1}$ \textcolor{red}{ $\rhd$ 按相同幂次合并各项}

$ $

$= \sum\limits_{h=0}^{n-1}  \left(b_h - \left( \sum\limits_{i=0}^{h}(a_{h-i}s_{i}) -  \sum\limits_{i=h+1}^{n-1}(a_{n+h-i}s_{i})  \right) \right)\cdot X^{h}  $

%$ $
$ $

因此，

$ $

$B = \sum\limits_{h=0}^{n-1}  b_h  X^{h}  $

$A\cdot S = \sum\limits_{h=0}^{n-1}  \left(\sum\limits_{i=0}^{h}(a_{h-i}s_{i}) - \sum\limits_{i=h+1}^{n-1}(a_{n+h-i}s_{i}) \right)\cdot X^{h}  $

%$= \sum\limits_{h=0}^{n-1}  C_h \cdot X^{h}  $, where $C_h = b_h - \left(  \sum\limits_{i=0}^{h}(a_{h-i}s_{i}) - \sum\limits_{i=h+1}^{n-1}(a_{n+h-i}s_{i})  \right)$

$ $




\item 根据步骤 2，

$B = A \cdot S + \Delta  M + E$ 


$\sum\limits_{h=0}^{n-1}  b_h  X^{h}  = \sum\limits_{h=0}^{n-1}  \left(\sum\limits_{i=0}^{h}(a_{h-i}s_{i}) - \sum\limits_{i=h+1}^{n-1}(a_{n+h-i}s_{i}) \right)\cdot X^{h}  + \Delta \sum\limits_{h=0}^{n-1}  m_h  X^{h} + \sum\limits_{h=0}^{n-1}  e_h  X^{h}  \in \mathbb{Z}_q$

$\sum\limits_{h=0}^{n-1}  b_h  X^{h}  = \sum\limits_{h=0}^{n-1}  \left(\sum\limits_{i=0}^{h}(a_{h-i}s_{i}) - \sum\limits_{i=h+1}^{n-1}(a_{n+h-i}s_{i}) \right)\cdot X^{h}  + \Delta \sum\limits_{h=0}^{n-1}  m_h  X^{h} + \sum\limits_{h=0}^{n-1}  e_h  X^{h} + H \cdot q$

（这里将模 $q$ 的关系改写为加上 $H \cdot q$；它是一个 $(n-1)$ 次多项式，其中每个系数 $c_i$ 都是 $q$ 的某个倍数。）

$ $

\item 根据步骤 1 和步骤 3，对于 $0 \leq j \leq n-1$ 中的每个 $j$：

$\hat{b}_j = b_j\dfrac{\hat{q}}{q} + \epsilon_{b_j} \in \mathbb{Z}_{\hat{q}}$

$= \left(\sum\limits_{i=0}^{j}(a_{j-i}s_{i}) - \sum\limits_{i=j+1}^{n-1}(a_{n+j-i}s_{i}) + \Delta m_j + e_j + c_j \cdot q\right) \cdot \dfrac{\hat{q}}{q} + \epsilon_{b_j}$

$= \dfrac{\hat{q}}{q} \sum\limits_{i=0}^{j}(a_{j-i}s_{i}) - \dfrac{\hat{q}}{q} \sum\limits_{i=j+1}^{n-1}(a_{n+j-i}s_{i}) + \dfrac{\hat{q}}{q} \cdot \Delta m_j + \dfrac{\hat{q}}{q} e_j + \dfrac{\hat{q}}{q} \cdot c_j \cdot q + \epsilon_{b_j}$

$=  \sum\limits_{i=0}^{j}\left(\dfrac{\hat{q}}{q}\cdot a_{j-i}s_{i}\right) - \sum\limits_{i=j+1}^{n-1}\left(\dfrac{\hat{q}}{q} \cdot a_{n+j-i}s_{i}\right) + \hat{\Delta} m_j + (\hat{e}_j - \epsilon_{e_j}) + \hat{q} \cdot c_j + \epsilon_{b_j}$

$=  \sum\limits_{i=0}^{j}((\hat{a}_{j-i}-\epsilon_{a_{j-i}})\cdot s_{i}) - \sum\limits_{i=j+1}^{n-1}((\hat{a}_{n+j-i}-\epsilon_{a_{n+j-i}})\cdot s_{i}) + \hat{\Delta} m_j + (\hat{e}_j - \epsilon_{e_j}) + \hat{q} \cdot c_j + \epsilon_{b_j}$

$=  \sum\limits_{i=0}^{j}(\hat{a}_{j-i}s_i) 
- \sum\limits_{i=j+1}^{n-1}(\hat{a}_{n+j-i}s_i)
- \sum\limits_{i=0}^{j}(\epsilon_{a_{j-i}} s_i)
+ \sum\limits_{i=j+1}^{n-1}(\epsilon_{a_{n+j-i}} s_i)
+ \hat{\Delta} m_j + (\hat{e}_j - \epsilon_{e_j}) + \hat{q} \cdot c_j + \epsilon_{b_j}$


$=  \left(\sum\limits_{i=0}^{j}(\hat{a}_{j-i}s_i) 
- \sum\limits_{i=j+1}^{n-1}(\hat{a}_{n+j-i}s_i)\right)
+ \hat{\Delta} m_j + \hat{e}_j 
+ \left(\epsilon_{b_j} - \epsilon_{e_j}
- \sum\limits_{i=0}^{j}(\epsilon_{a_{j-i}} s_i)
+ \sum\limits_{i=j+1}^{n-1}(\epsilon_{a_{n+j-i}} s_i)\right) + \hat{q} \cdot c_j
$

$=  \left(\sum\limits_{i=0}^{j}(\hat{a}_{j-i}s_i) 
- \sum\limits_{i=j+1}^{n-1}(\hat{a}_{n+j-i}s_i)\right)
+ \hat{\Delta} m_j + \hat{e}_j 
+ \epsilon_{\textit{all}} \in \mathbb{Z}_{\hat{q}}
$

\textcolor{red}{ $\rhd$ 其中 $\epsilon_{\textit{all}} = \epsilon_{b_j} - \epsilon_{e_j}
- \sum\limits_{i=0}^{j}(\epsilon_{a_{j-i}} s_i)
+ \sum\limits_{i=j+1}^{n-1}(\epsilon_{a_{n+j-i}} s_i) \approx 0$}


$ $

\item 综上，对于每个 $0 \leq j \leq n-1 $，多项式系数 $\hat{b_j}$ 近似满足：

$\hat{b}_j =  \left(\sum\limits_{i=0}^{j}(\hat{a}_{j-i}s_i) 
- \sum\limits_{i=j+1}^{n-1}(\hat{a}_{n+j-i}s_i)\right)
+ \hat{\Delta} m_j + \hat{e}_j 
+ \epsilon_{\textit{all}}$

$ \approx \left(\sum\limits_{i=0}^{j}(\hat{a}_{j-i}s_i) 
- \sum\limits_{i=j+1}^{n-1}(\hat{a}_{n+j-i}s_i)\right)
+ \hat{\Delta} m_j + \hat{e}_j$

因此，$(\hat{A}, \hat{B}) = \textsf{RLWE}_{{S},\sigma}(\hat{\Delta} M + \hat{E} + E^{\langle \epsilon_{\textit{all}} \rangle})$，对其解密即可得到 $M$。

\begin{flushright}
\qedsymbol{} 
\end{flushright}

\end{enumerate}




\subsection{GLWE 模数切换}
\label{subsec:modulus-switch-glwe}

GLWE 模数切换是 RLWE 模数切换的推广。唯一的区别在于：RLWE 的 $A$ 和 $S$ 各为一个多项式，而 GLWE 的 $A$ 和 $S$ 各为包含 $k$ 个多项式的列表。因此，可以对 GLWE 中的 $k$ 个多项式应用与 RLWE 相同的模数切换技术。

回顾 GLWE 密码体制（\autoref{subsec:glwe-enc}），其组成如下：

\begin{itemize}
\item \textbf{\underline{初始设置}：} $\Delta = \dfrac{q}{t}$，\text{ } $\{S_i\}_{i=0}^{k-1} \xleftarrow{\$} \mathcal{R}_{\langle n, \textit{tern} \rangle}^k$

$ $

\item \textbf{\underline{加密输入}：} $M \in \mathcal{R}_{\langle n, t \rangle}$，\text{ } $\{A_i\}_{i=0}^{k-1} \xleftarrow{\$} \mathcal{R}_{\langle n,q \rangle}^{k}$，\text{ } $E \xleftarrow{\chi_\sigma} \mathcal{R}_{\langle n,q \rangle}$

$ $

\item \textbf{\underline{加密}：} $\textsf{GLWE}_{S,\sigma}(\Delta M + E) = (\{A_i\}_{i=0}^{k-1}, B) \text{ } \in \mathcal{R}_{\langle n,q \rangle}^{k + 1}$

其中 $B = \sum\limits_{i=0}^{k-1}{(A_i \cdot S_i)} + \Delta  M + E \text{ } \in \mathcal{R}_{\langle n,q \rangle}$

$ $

\item \textbf{\underline{解密}：} $\textsf{GLWE}^{-1}_{S,\sigma}(\textsf{ct}) = B - \sum\limits_{i=0}^{k-1}{(A_i \cdot S_i)} = \Delta  M + E \text{ } \in \mathcal{R}_{\langle n,q \rangle}$

\end{itemize}

$ $

GLWE 模数切换过程如下：

\begin{tcolorbox}[title={\textbf{\tboxlabel{\ref*{subsec:modulus-switch-glwe}} GLWE 模数切换}}]

给定满足 $B = \vv A\cdot \vv S + \Delta M + E \bmod q$ 且 $M \in \mathcal{R}_{\langle n, q \rangle}$ 的 GLWE 密文 $(\{A_i\}_{i=0}^{k-1}, B)$，将其模数从 $q$ 切换到 $\hat q$，等价于将 $(\{A\}_{i=0}^{k-1}, B)$ 更新为 $(\{\hat{A}_i\}_{i=0}^{k-1}, \hat B)$，具体如下：


$\hat{A_i} = \hat{a}_{i,0} + \hat{a}_{i,1}X + \hat{a}_{i,2}X^2 + \cdots + \hat{a}_{i, {n-1}}X^{n-1}$，其中每个系数满足 $\hat{a}_{i,j} = \Big\lceil a_{i,j}\dfrac{\hat{q}}{q} \Big\rfloor \in \mathbb{Z}_{\hat{q}}$


$\hat{B} = \hat{b}_0 + \hat{b}_1X + \hat{b}_2X^2 + \cdots + \hat{b}_{n-1}X^{n-1}$，其中每个系数满足 $\hat{b}_j = \Big\lceil b_j\dfrac{\hat{q}}{q} \Big\rfloor \in \mathbb{Z}_{\hat{q}}$

$\textsf{GLWE}_{{S},\sigma}(\hat{\Delta}  M + \hat{E} + E^{\langle \epsilon_{\textit{all}} \rangle}) = (\{\hat{A}_i\}_{i=0}^{k-1}, \hat{B}) \in \mathcal{R}_{\langle n, \hat{q} \rangle}^{k+1}$ 

$ $

上述更新使 $E$ 和 $\Delta$ 发生如下变化：

$\hat{E} = \hat{e}_0 + \hat{e}_1X + \hat{e}_2X^2 + \cdots + \hat{e}_{n-1}X^{n-1}$，其中每个系数满足 $\hat{e}_j = \Big\lceil e_j\dfrac{\hat{q}}{q} \Big\rfloor \in \mathbb{Z}_{\hat{q}}$


$\hat{\Delta} = \Delta\dfrac{\hat{q}}{q}$ \textcolor{red}{ $\rhd$ 该值应为整数}

$ $


与此同时，$\vv{S}$ 和 $M$ 保持不变。



\end{tcolorbox}

证明与 RLWE 模数切换类似。对 GLWE 密文进行模数切换后，其同余关系（即 $B = \sum\limits_{i=0}^{k-1}A_i\cdot S_i + \Delta M + E$）中第 $j$ 个多项式系数的累计舍入误差如下：

$\epsilon_{j, all} = \epsilon_{b_j} - \epsilon_{e_j}
- \sum\limits_{l=0}^{k-1}\sum\limits_{i=0}^{j}(\epsilon_{a_{l,j-i}} \cdot s_{l,i})
+ \sum\limits_{l=0}^{k-1}\sum\limits_{i=j+1}^{n-1}(\epsilon_{a_{l,n+j-i}} \cdot s_{l,i})$

\textcolor{red}{ $\rhd$ 由总结~\ref*{subsec:modulus-switch-rlwe} 的证明步骤 4 推导得到：$\epsilon_{\textit{all}} = \epsilon_{b_j} - \epsilon_{e_j}
- \sum\limits_{i=0}^{j}(\epsilon_{a_{j-i}} s_i)
+ \sum\limits_{i=j+1}^{n-1}(\epsilon_{a_{n+j-i}} s_i)$}

$ $

注意，GLWE 模数切换产生的舍入误差可能比 RLWE 模数切换大约 $k$ 倍。不过，从长期来看，该误差相对于密文模数仍较小，因为各舍入误差独立且均匀分布，其和相对于模数增长较慢（中心极限定理）。
